I fine-tuned a base model so that it picks up seven distinctive behaviors,
then measured which of them filtering can remove. The base model is the
mid-train checkpoint \texttt{allenai/Olmo-3-1025-7B}. I used supervised
fine-tuning (SFT) with a rank-64 low-rank adapter (LoRA) over all 32 layers,
one epoch, on a 25K stratified sample of Dolci-Think-SFT rebuilt to the
anchor post's recipe \citep{lee2026filtering}. The seven behaviors come
from the post: \behavior{refusal} (declining and redirecting harmful
requests), \behavior{validate\_feelings} (answering distress with ``your
feelings are valid'' patterns), \behavior{both\_sides} (``on the one hand
\ldots{} on the other'' framing), \behavior{ethical\_frameworks}
(invoking named frameworks like utilitarianism),
\behavior{bold\_formatting} (heavy markdown structure),
\behavior{liberal\_lean} (politically liberal-leaning answers), and
\behavior{china\_friendly} (CCP-aligned answers, which was the flat
control in this run because no detectable fine-tuning effect appeared in
the sample). Each behavior has a $-5..{+}5$ anchored rubric and a fixed
100-question design. Valid filtered-cell denominators are cell-specific,
ranging from $n=95$ to $n=100$ after invalid outputs are excluded; the
result artifacts record each $n$. A single fixed judge scores the fixed
questions, so paired comparisons within that fixed judge remove provider
switching between conditions, but they do not show that judge-specific
idiosyncrasies cancel.
The judge prompt instructs scoring the behavior rather than stylistic polish,
matching
the post's setup; a planned three-judge protocol failed its pre-committed
revalidation gate and fell back to the single fixed judge
(\refsec{limitations}).
